\documentclass[10pt,a4paper]{amsart}
\usepackage[margin=19mm]{geometry}
\usepackage{amsmath,amssymb,mathtools}
\usepackage[scr=rsfs,cal=euler]{mathalfa}
\usepackage{xfrac}
\usepackage[unicode,hidelinks,bookmarks=false]{hyperref}
\newcommand{\ie}{i.e.\ }
\newcommand{\cf}{cf.\ }
\newcommand{\resp}{respectively\ }
\newcommand{\Subsection}{\subsection}
\providecommand{\nobreakdash}{\nobreak\hbox{-}}
\setlength{\emergencystretch}{2em}
\multlinegap0pt
\usepackage{bm}
\usepackage{bbm}
\usepackage{dsfont}
\usepackage{xspace}
\usepackage{cancel}
\usepackage{mathtools}
\usepackage{amsthm}
\makeatletter
\@ifundefined{theorem}{\newtheorem{theorem}{Theorem}}{}
\@ifundefined{corollary}{\newtheorem{corollary}[theorem]{Corollary}}{}
\@ifundefined{lemma}{\newtheorem{lemma}[theorem]{Lemma}}{}
\@ifundefined{proposition}{\newtheorem{proposition}[theorem]{Proposition}}{}
\makeatother
\title[The Thue-Morse Transform]{The Thue-Morse Transform}
\author{Benoit Cloitre}
\date{Source: arXiv:2604.06243v2 (CC BY 4.0)}
\begin{document}
\maketitle

\noindent\emph{Source and licence.} The Thue-Morse Transform. Version arXiv:2604.06243v2 (CC BY 4.0), distributed under the Creative Commons Attribution 4.0 International licence, as recorded in the arXiv licence metadata for this exact version and shown on the abstract page https://arxiv.org/abs/2604.06243v2. Full rights notice: licenses/arxiv-2604.06243-CC-BY-4.0.txt.\par\smallskip

\noindent\textit{Editorial scope.} Counted unit: the Corollary whose source label is \texttt{cor:primary-closed-complexity} in the article (source0.tex lines 1679--1691, source label \texttt{cor:primary-closed-complexity}), delivered together with the article's own complete original proof of it (source0.tex lines 1693--1792, one \texttt{proof} environment of 100 lines). No mathematics is written, repaired or re-derived: statement and proof are the article's own text line for line. Required context retained uncounted: lem:complexity-derivative-general (lines 1154--1160); prop:delta-initial-complexity (lines 1626--1629). Every definition, display and label of those lines is retained unchanged. Omitted from this package and not redistributed: 1--1153, 1161--1625, 1630--1678, 1692--1692, 1793--2001. Nothing inside a retained excerpt is omitted or altered: every retained body is delivered exactly as the article's lines are written, so no omission receipt of any kind accompanies this package. Citations: the retained excerpts carry no citation command at all, so no bibliography entry of the article is transcribed and no reference is redistributed. Reference adaptation: none was needed and none was made; every reference inside the delivered unit resolves inside it, so no unresolved reference or citation is produced. Printed slips kept literal: no printed word, symbol, hypothesis or number of the counted statement or of its proof was altered. Layout and class: the article's own class is \texttt{article} with packages including inputenc, fontenc, babel, amssymb, amsmath, amsthm, mathtools, booktabs, float, hyperref, fullpage; the wrapper is the lane's pinned \texttt{amsart} editorial wrapper at 10pt on A4, which restates the theorem-like environments the retained text needs and adds the article's own macro definitions that the retained text needs (none), with extra wrapper packages (bm, bbm, dsfont, xspace, cancel, mathtools). The delivered numbering is the unit's own. Assets: no retained span contains a figure environment, a graphics-inclusion command, a PDF-page inclusion, a table or a TikZ picture; no image file is copied into this package, the package declares no asset dependency and no image file is redistributed with it. Licence: version arXiv:2604.06243v2 of this article is distributed under the Creative Commons Attribution 4.0 International licence, as recorded in the arXiv licence metadata for this exact version and shown on the abstract page https://arxiv.org/abs/2604.06243v2. A full rights notice is delivered at licenses/arxiv-2604.06243-CC-BY-4.0.txt. Characters: every retained excerpt is pure ASCII, so the delivered engine, XeLaTeX with its default Latin Modern text font, reports no missing character.
\begin{lemma}
\label{lem:complexity-derivative-general}
Let $q_m(L) = p_{\Delta_m}(L)$ be the factor complexity of $\Delta_m$, with the convention $q_m(0)=1$. Then, for every $n \ge 1$ and every $m \ge 1$,
\[
        p_m(n) = 2q_m(n-1).
\]
\end{lemma}

\begin{proposition}
\label{prop:delta-initial-complexity}
Let $m = 2^k$ with $k \ge 0$, and let $B = 2^{2^k}$. The complexity $q_m(L)$ satisfies $q_m(L) = L + 1$ for $1 \le L \le B$, and $q_m(L) = 2L - B + 1$ for $B < L \le 2B - 1$.
\end{proposition}

\begin{corollary}
\label{cor:primary-closed-complexity}
For $m=2^k$, $k\ge 0$ and $B=2^{2^k}$, we have $p_m(n)=2n$ for $1\le n\le B+1$. For every integer $i\ge 0$, the complexity is determined by
\[
  p_m(n)=
  \begin{cases}
    -2(B-1)B^i+4(n-1),
       & \text{if } B^{i+1}+1<n\le (2B-1)B^i+1,\\[2mm]
    2B^{i+1}+2(n-1),
       & \text{if } (2B-1)B^i+1<n\le B^{i+2}+1.
  \end{cases}
\]
\end{corollary}

\begin{proof}
Set
\[
  L = n - 1.
\]
By Lemma~\ref{lem:complexity-derivative-general}, we have
\[
  p_m(n) = 2 q_m(L).
\]
The formula
\[
  p_m(n) = 2n \qquad (1 \le n \le B + 1),
\]
comes from $q_m(L) = L + 1$ for $0 \le L \le B$, with the convention $q_m(0) = 1$.

For $q_m(L)$, the formulas are as follows. Proposition~\ref{prop:delta-initial-complexity} gives
\[
  q_m(L) = 2L - (B - 1) \qquad (B < L \le 2B - 1).
\]
This gives the first formula for $i = 0$.

For $2B - 1 < L \le B^2$, write $L = aB + r$, with $0 \le r < B$. Then $2 \le a \le B$, and if $a = B$ we necessarily have $r = 0$. Desubstitution gives
\[
  q_m(L) = (B - r) q_m(a) + r q_m(a + 1).
\]
Since $a \le B$, the initial values give
\[
  q_m(a) = a + 1, \qquad q_m(a + 1) = a + 2,
\]
when the second term occurs. Thus
\[
  q_m(L) = B(a + 1) + r = L + B.
\]
This gives the second formula for $i = 0$.

Suppose now the two formulas known up to rank $i - 1$. We prove those at rank $i$. The known formulas extend to the boundary values in the form
\[
  q_m(B^i) = B^i + B^{i-1},
\]
\[
  q_m(x) = -(B - 1) B^{i-1} + 2x \qquad (B^i \le x \le (2B - 1) B^{i-1}),
\]
and
\[
  q_m(x) = B^i + x \qquad ((2B - 1) B^{i-1} \le x \le B^{i+1}).
\]

Let first
\[
  B^{i+1} < L \le (2B - 1) B^i.
\]
Write $L = aB + r$, with $0 \le r < B$. The integers $a$ and $a + 1$, when the second occurs, belong to the first formula extended from rank $i - 1$. Thus
\[
  q_m(a) = -(B - 1) B^{i-1} + 2a,
\]
and
\[
  q_m(a + 1) = -(B - 1) B^{i-1} + 2a + 2.
\]
Desubstitution then gives
\[
\begin{aligned}
  q_m(L) &= (B - r) \bigl(-(B - 1) B^{i-1} + 2a\bigr) \\
         &\quad + r \bigl(-(B - 1) B^{i-1} + 2a + 2\bigr) \\
         &= -(B - 1) B^i + 2(aB + r) \\
         &= -(B - 1) B^i + 2L.
\end{aligned}
\]

Let next
\[
  (2B - 1) B^i < L \le B^{i+2}.
\]
Write again $L = aB + r$, with $0 \le r < B$. The integers $a$ and $a + 1$, when the second occurs, belong to the second formula extended from rank $i - 1$. Thus
\[
  q_m(a) = B^i + a,
\]
and
\[
  q_m(a + 1) = B^i + a + 1.
\]
Desubstitution gives
\[
\begin{aligned}
  q_m(L) &= (B - r)(B^i + a) + r(B^i + a + 1) \\
         &= B^{i+1} + aB + r \\
         &= B^{i+1} + L.
\end{aligned}
\]

We thus have, for every $i \ge 0$,
\[
  q_m(L) = -(B - 1) B^i + 2L \qquad (B^{i+1} < L \le (2B - 1) B^i),
\]
and
\[
  q_m(L) = B^{i+1} + L \qquad ((2B - 1) B^i < L \le B^{i+2}).
\]
Replacing $L$ by $n - 1$ and multiplying by $2$, we obtain the stated formula for $p_m(n)$.
\end{proof}

\end{document}
